\subsection{无迹极化与自旋关联的张量恢复}

零极化多项式隐藏了迹项，二点函数和守恒条件却会重新调用这些项。为使张量缩并自足，固定恢复规则。

\begin{definition}[调和投影与零锥导数]
令 $F_S(z)$ 为 $d$ 维辅助向量 $z$ 的齐次 $S$ 次多项式。其对称无迹投影为
\begin{equation}\label{eq:ghs-r4-stf-projector}
 \Pi_SF_S=\sum_{k=0}^{\lfloor S/2\rfloor}
 \frac{(-1)^k\Gamma(S+d/2-k-1)}{4^kk!\Gamma(S+d/2-1)}
 (z^2)^k(\partial_z^2)^kF_S.
\end{equation}
$S=0$ 单独定义 $\Pi_0=1$。令 $\nu=d/2-1$，作用于零锥等价类的 Thomas 导数为
\[
 \mathcal D_{z,\mu}=(\nu+z\cdot\partial_z)\partial_{z,\mu}
 -\tfrac12z_\mu\partial_z^2.
\]
它将 $z^2$ 倍数映为 $z^2$ 倍数，因此 $\partial_x^\mu\mathcal D_{z,\mu}J_S$ 与多项式代表无关。完整投影递推与任意秩证明可参照《群论》的对称无迹张量部分。
\end{definition}
\begin{proof}
对齐次 $m$ 次多项式 $G_m$，
\[
 \partial_z^2[(z^2)^kG_m]=(z^2)^k\partial_z^2G_m
 +4k(m+k+d/2-1)(z^2)^{k-1}G_m.
\]
代入 $m=S-2k$，相邻迹阶系数满足
$a_{k+1}=-a_k/[4(k+1)(S-k+d/2-2)]$，$a_0=1$，得到 $\Gamma$ 形式；最高阶以后重复 Laplace 运算为零。

对 Thomas 导数展开 $\mathcal D_{z,\mu}(z^2G_m)$。不含 $z^2$ 的项为
$[2(\nu+m+1)-(d+2m)]z_\mu G_m=0$，所以改变零锥多项式代表不会改变散度条件。
\end{proof}
